Các giải pháp hiện tại cho quan sát lớp học có thể chia thành bốn nhóm chính: quan sát thủ công, nhận diện hành vi bằng thị giác máy tính, khám phá cấu trúc trên chuỗi thời gian và sử dụng mô hình ngôn ngữ lớn để hỗ trợ suy luận đồ thị. Mỗi nhóm giải quyết một phần của bài toán, nhưng chưa tạo thành một quy trình thống nhất từ video thô đến báo cáo quan sát có thể kiểm chứng.

Nhóm giải pháp truyền thống dựa trên người quan sát được huấn luyện và các bộ mã hành vi định nghĩa trước. Ưu điểm của nhóm này là có khả năng gắn quan sát với bối cảnh sư phạm và tiêu chí nghiên cứu rõ ràng. Tuy nhiên, hạn chế lớn là chi phí nhân lực cao, thời gian xử lý dài và khó duy trì độ nhất quán khi mở rộng sang nhiều lớp học hoặc nhiều người chấm \,\cite{volpe2025classroom}. Vì vậy, cách tiếp cận này phù hợp với phân tích chuyên sâu quy mô nhỏ hơn là xử lý lượng lớn video lớp học.

Nhóm giải pháp thứ hai dùng thị giác máy tính để nhận diện hành vi học sinh trong video. Các nghiên cứu gần đây cho thấy detector dựa trên YOLO và các biến thể có attention có thể cải thiện khả năng phát hiện trong bối cảnh lớp học đông người, nhiều che khuất và nhiều hành vi có tín hiệu thị giác gần giống nhau \,\cite{wang2023smartclassroom,yang2023yolov7bra,chen2023improvedyolov8,zhang2025yolov8attention}. Dù vậy, phần lớn hệ thống loại này dừng ở nhãn, hộp phát hiện, số đếm hoặc thống kê tần suất. Các đầu ra này hữu ích nhưng chưa trả lời đầy đủ câu hỏi hành vi nào thường đi trước hành vi nào, hoặc các tín hiệu hành vi tạo thành cấu trúc thời gian ra sao.

Nhóm giải pháp thứ ba là các phương pháp mô hình hóa chuỗi thời gian và khám phá liên kết dự báo. Các phương pháp theo tinh thần Granger hoặc các mô hình đồ thị thời gian có thể phát hiện việc lịch sử của một biến giúp dự báo biến khác \,\cite{granger1969causality,runge2023causalreview}. Tuy nhiên, các phương pháp này thường giả định đầu vào đã là chuỗi đa biến sạch và có ý nghĩa, trong khi dữ liệu video lớp học phải đi qua các bước phát hiện, theo dõi, gom sự kiện và chuẩn hóa. Nếu thiếu cơ chế truy vết về video gốc, đồ thị thu được có thể khó kiểm tra và dễ bị diễn giải tách khỏi bằng chứng thị giác.

Nhóm giải pháp thứ tư dùng mô hình ngôn ngữ lớn như nguồn tri thức miền hoặc bộ đề xuất cấu trúc. LLM có thể giúp nêu giả thuyết về quan hệ hợp lý giữa các biến hành vi, giảm không gian tìm kiếm hoặc giải thích đồ thị bằng ngôn ngữ tự nhiên \,\cite{darvariu2024llmpriors,roy2025causalllm}. Hạn chế của hướng này là đầu ra của LLM có thể thiên kiến, thiếu nhất quán hoặc đề xuất quan hệ nghe hợp lý nhưng không được dữ liệu ủng hộ. Do đó, LLM không nên được dùng như nguồn quyết định đồ thị cuối cùng nếu thiếu kiểm chứng định lượng.

Từ các hạn chế trên, có thể thấy khoảng trống chính không chỉ nằm ở việc xây dựng một detector chính xác hơn. Bài toán cần một quy trình nối liền bốn tầng: phát hiện hành vi có thể quan sát trong video, chuyển phát hiện thành chuỗi thời gian có thể truy vết, khám phá đồ thị liên kết dự báo bằng mô hình định lượng và trình bày kết quả bằng báo cáo thận trọng. CausClass được đề xuất để đáp ứng khoảng trống này bằng cách kết hợp YOLO26n-MHSA, ByteTrack, biểu diễn hành vi theo time bin, AERCA có mask, tìm kiếm beam có hướng dẫn bởi LLM và cơ chế sinh báo cáo theo nguyên tắc nêu bằng chứng trước.